This thesis systematically analyzes the key design factors of geometry-model-based feature distillation, which enhances the 3D spatial perception of Multimodal Large Language Models (MLLMs) without incurring additional memory or computational overhead at inference time. Across the R2R-CE and RxR-CE benchmarks, we evaluate how teacher model selection, information flow directionality, and the temporal stride of navigation history affect Vision-and-Language Navigation (VLN) performance. Our findings demonstrate that: (1) alignment between the teacher model's geometry decoding architecture and the distillation target layer strongly impacts performance; (2) the information flow directionality mismatch between bidirectional and causal teachers is not a dominant bottleneck in the current pipeline; and (3) a wider context stride helps the MLLM capture longer-horizon navigation context, but reduces inter-frame visual overlap and diminishes 3D spatial reasoning gains, creating a structural trade-off.

Future work should more rigorously validate these factors and incorporate them into improved distillation designs. Specifically, multi-layer distillation, which incorporates intermediate layers into the distillation target when using DPT-head-based teachers, and a decoupled stride design that separates the MLLM's long-horizon context retention from the geometry teacher's high-quality 3D representation generation represent promising directions for improvement.
